\documentclass[11pt]{amsart}

\usepackage[T1]{fontenc}
\usepackage{lmodern}
\usepackage{microtype}
\usepackage{amsmath,amssymb}
\usepackage{booktabs}
\usepackage[hidelinks]{hyperref}

\newtheorem{theorem}{Theorem}
\newtheorem{proposition}[theorem]{Proposition}
\theoremstyle{remark}
\newtheorem{remark}[theorem]{Remark}

\newcommand{\W}{\mathcal W}

\title[A seven-vertex counterexample]
{A seven-vertex counterexample to Conjecture~5.3 of G\"obel and Misra}

\date{}

\subjclass[2020]{05C38, 15A15}
\keywords{colored path, tridiagonal determinant, matching expansion,
counterexample}

\begin{document}

\begin{abstract}
G\"obel and Misra conjectured that every equality between determinant
polynomials of colored paths is generated by identity, reflection, or one of
two parity-dependent constructions. We refute the odd-order clause of their
Conjecture~5.3. Two uniformly vertex-colored paths on seven vertices, with
edge words
\[
  bbcbcc \qquad\text{and}\qquad bccbbc,
\]
have the same determinant polynomial but satisfy none of the conjectured
alternatives. An exact exhaustive computation shows that seven vertices are
minimal. The even-order clause remains open.
\end{abstract}

\maketitle

Let a colored path on $m$ vertices have vertex word
$u=u_1\cdots u_m$ and edge word $w=w_1\cdots w_{m-1}$, with vertex and
edge colors drawn from disjoint alphabets. Associate variables $x_c$ to
vertex colors and $y_d$ to edge colors, and write
\[
K(u,w)=
\begin{pmatrix}
 x_{u_1} & y_{w_1} \\
 y_{w_1} & x_{u_2} & \ddots \\
 & \ddots & \ddots & y_{w_{m-1}} \\
 & & y_{w_{m-1}} & x_{u_m}
\end{pmatrix},
\qquad
D(u,w)=\det K(u,w).
\]
If $M$ is a matching of the path and $V(M)$ is the set of vertices incident
with $M$, then
\begin{equation}\label{eq:matching}
D(u,w)=
\sum_M(-1)^{|M|}
\prod_{i\in M}y_{w_i}^{2}
\prod_{j\notin V(M)}x_{u_j}.
\end{equation}
This is the matching expansion in~\cite[Lemma~3.3]{GM}.

Conjecture~5.3(2) of~\cite{GM} asserts that, for odd $m$, equality of two
such determinant polynomials arises only from identity, reflection, or the
coloring configuration of~\cite[Theorem~3.8]{GM}, up to reflection. In
reading that theorem, we replace the edge-index ranges in its condition~(3)
by $\{1,\ldots,m-1\}$, as required for a path on $m$ vertices and as used by
the matrices in its proof. In particular, its edge condition makes each
parity class of edge positions monochromatic and exchanges the two colors
between the paths.

\begin{theorem}\label{thm:counterexample}
Let $a$ be a vertex color and let $b\ne c$ be edge colors. Put
\[
 p=bbcbcc,
 \qquad
 q=bccbbc,
 \qquad
 x=x_a,
 \qquad
 B=y_b^2,
 \qquad
 C=y_c^2.
\]
Then
\begin{equation}\label{eq:identity}
D(a^7,p)=D(a^7,q)
=x^7-3x^5(B+C)+2x^3(B^2+3BC+C^2)-2xBC(B+C).
\end{equation}
The two colored paths satisfy none of the alternatives in
Conjecture~5.3(2). Hence that clause, and therefore Conjecture~5.3 as stated,
is false.
\end{theorem}

\begin{proof}
For a six-letter edge word $z$, define
\[
 A_r(z)=
 \sum_{\substack{M\text{ a matching}\\ |M|=r}}
 \prod_{i\in M}y_{z_i}^{2}.
\]
Since all seven vertices have the same color,~\eqref{eq:matching} gives
\[
 D(a^7,z)=x^7-x^5A_1(z)+x^3A_2(z)-xA_3(z).
\]
The size-two matchings are indexed by
\[
 13,14,15,16,24,25,26,35,36,46,
\]
and the size-three matchings by
\[
 135,136,146,246.
\]
Substitution of either $p$ or $q$ therefore gives
\[
\begin{aligned}
 A_1(p)=A_1(q)&=3B+3C,\\
 A_2(p)=A_2(q)&=2B^2+6BC+2C^2,\\
 A_3(p)=A_3(q)&=2B^2C+2BC^2.
\end{aligned}
\]
This proves~\eqref{eq:identity}.

The edge words are distinct, and the reverse of $p$ is $ccbcbb$, not $q$.
Their odd- and even-position edge subsequences are
\[
\begin{array}{c|cc}
 & \text{odd positions} & \text{even positions}\\
\hline
p & (b,c,c) & (b,b,c)\\
q & (b,c,b) & (c,b,c).
\end{array}
\]
None is monochromatic. Reversal only reverses and exchanges these two
subsequences, so no choice of orientations, and no interchange of the two
paths, satisfies the edge condition of Theorem~3.8. Thus none of the
conjectured alternatives applies.
\end{proof}

\begin{remark}\label{rem:pd}
For every real specialization with
$x>2\max\{|y_b|,|y_c|\}$, both matrices in
Theorem~\ref{thm:counterexample} are positive definite by strict diagonal
dominance. Thus the equality occurs inside the positive-definite parameter
space of the corresponding colored Gaussian graphical models.
\end{remark}

\section*{Minimality}

For $n\geq1$, let $\W_n$ be the set of all surjective words on
$\{1,\ldots,r\}$, for some $1\leq r\leq n$, whose color multiplicities are
weakly decreasing. When multiplicities tie, all label orders are retained.

\begin{proposition}\label{prop:minimal}
No counterexample to Conjecture~5.3 has at most six vertices, under the
orientation-invariant reading of the two constructions used above.
Consequently, the counterexample in Theorem~\ref{thm:counterexample} is
smallest with respect to the number of vertices.
\end{proposition}

\begin{proof}
Equality of determinant polynomials forces equality of the vertex-color
multiplicities, because the edge-degree-zero term in~\eqref{eq:matching} is
$\prod_j x_{u_j}$. It also forces equality of the edge-color multiplicities:
after setting every vertex variable equal to $1$, the coefficient of the pure
monomial $y_d^2$ is minus the number of edges of color $d$. A common
relabeling of the vertex colors and, independently, of the edge colors
therefore sends every candidate pair to
$\W_m\times\W_{m-1}$.

For each $(u,w)\in\W_m\times\W_{m-1}$ with $3\leq m\leq6$, we evaluated
\eqref{eq:matching} as an exact integer coefficient dictionary indexed by the
full exponent vector and grouped paths with identical dictionaries. Every
unordered pair of distinct paths in a group was first tested for reflection.
Each remaining pair was then tested against Theorem~3.6 when $m$ is even and
the type-correct conditions of Theorem~3.8 when $m$ is odd, allowing the two
paths to be interchanged and each path to be independently reversed. The
disjoint counts are
\begin{center}
\small
\setlength{\tabcolsep}{5pt}
\begin{tabular}{crrrrr}
\toprule
$m$ & paths & equal pairs & reflections & constructions & residual \\
\midrule
3 & 30        & 14        & 14        & 0   & 0 \\
4 & 470       & 248       & 232       & 16  & 0 \\
5 & 11{,}562  & 5{,}780   & 5{,}772   & 8   & 0 \\
6 & 394{,}092 & 197{,}496 & 197{,}016 & 480 & 0 \\
\bottomrule
\end{tabular}
\end{center}
The cases $m\leq2$ are immediate from
$D(u_1,\varnothing)=x_{u_1}$ and
$D(u_1u_2,w_1)=x_{u_1}x_{u_2}-y_{w_1}^2$.

\end{proof}

The counterexample refutes precisely the odd-order clause of
Conjecture~5.3. The exhaustive computation verifies the even-order clause
only through six vertices; the general even-order clause remains open.

\begin{thebibliography}{9}

\bibitem{GM}
H.~G\"obel and P.~Misra,
\emph{Linear relations of colored Gaussian cycles},
arXiv:2506.23936v1 [math.CO], 2025.

\end{thebibliography}

\end{document}